% ------------------------------------------------------------------------
% file `2_noel_2021_part2-exercise-16-exercise-body.tex'
%   in folder `xsim/'
%
%     exercise of type `exercise' with id `16'
%
% generated by the `exercise' environment of the
%   `xsim' package v0.21 (2022/02/12)
% from source `2_noel_2021_part2' on 2022/11/30 on line 68
% ------------------------------------------------------------------------
    À l'aide du tableau ci-dessous, réponds aux questions.
    \begin{center}
        \begin{tblr}{hlines, vlines,row{1}={m,font=\bfseries}, columns={c, 4cm}, rowsep=1ex}
            Planète & Distance par rapport au soleil (km) & Notation scientifique\\
            %Mercure & \num{58000000}       \\
            Vénus   & \num{108 190 000} &   \\
            Terre   & \num{149 569 000} &   \\
            Mars    & \num{227 940 000} &   \\
            Jupiter & \num{778 000 000} &   \\
            %Saturne & \num{1 427 000 000}  \\
            %Uranus  & \num{2 871 000 000}  \\
            %Neptune & \num{4 497 000 000}
        \end{tblr}
    \end{center}
    \begin{enumerate}[itemsep=3em]
        \item Convertis ces distances en notation scientifique.
        \item Laquelle de ces planètes se trouve à environ $2\cdot 10^8$ km du soleil ?
        \item Dans le système solaire, on peut retrouver des milliers d'astéroïdes. Ceux-ci sont beaucoup plus petits que les planètes. Il y a une immense ceinture d'astéroïdes à environ $6\cdot 10^8$ km du soleil. Entre quelles planètes se trouvent cette ceinture ?
    \end{enumerate}
